\documentclass[10pt,a4paper]{amsart}
\usepackage{amsmath,amssymb,amsthm,graphicx}
% Retained minimal preamble for the Stacks Project chapter varieties.tex.
% Verbatim excerpt of the chapter preamble: Xy-pic setup (native preamble lines 14-18), the native theorem
% declarations with the mathematical macros (native preamble lines 154-217), and the project own no-op
% definitions of the reference/slogan/history display environments (native wrapper.tex lines 1-3) that the
% counted statement uses for its slogan. Copied mechanically; nothing added, removed or redefined.
% For commutative diagrams we use Xy-pic
\usepackage[all]{xy}

% We use 2cell for 2-commutative diagrams.
\xyoption{2cell}
% Theorem environments.
%
\theoremstyle{plain}
\newtheorem{theorem}[subsection]{Theorem}
\newtheorem{proposition}[subsection]{Proposition}
\newtheorem{lemma}[subsection]{Lemma}

\theoremstyle{definition}
\newtheorem{definition}[subsection]{Definition}
\newtheorem{example}[subsection]{Example}
\newtheorem{exercise}[subsection]{Exercise}
\newtheorem{situation}[subsection]{Situation}

\theoremstyle{remark}
\newtheorem{remark}[subsection]{Remark}
\newtheorem{remarks}[subsection]{Remarks}

\numberwithin{equation}{subsection}

% Macros
%
\def\lim{\mathop{\mathrm{lim}}\nolimits}
\def\colim{\mathop{\mathrm{colim}}\nolimits}
\def\Spec{\mathop{\mathrm{Spec}}}
\def\Hom{\mathop{\mathrm{Hom}}\nolimits}
\def\Ext{\mathop{\mathrm{Ext}}\nolimits}
\def\SheafHom{\mathop{\mathcal{H}\!\mathit{om}}\nolimits}
\def\SheafExt{\mathop{\mathcal{E}\!\mathit{xt}}\nolimits}
\def\Sch{\mathit{Sch}}
\def\Mor{\mathop{\mathrm{Mor}}\nolimits}
\def\Ob{\mathop{\mathrm{Ob}}\nolimits}
\def\Sh{\mathop{\mathit{Sh}}\nolimits}
\def\NL{\mathop{N\!L}\nolimits}
\def\CH{\mathop{\mathrm{CH}}\nolimits}
\def\proetale{{pro\text{-}\acute{e}tale}}
\def\etale{{\acute{e}tale}}
\def\QCoh{\mathit{QCoh}}
\def\Ker{\mathop{\mathrm{Ker}}}
\def\Im{\mathop{\mathrm{Im}}}
\def\Coker{\mathop{\mathrm{Coker}}}
\def\Coim{\mathop{\mathrm{Coim}}}

% Boxtimes
%
\DeclareMathSymbol{\boxtimes}{\mathbin}{AMSa}{"02}

%
% Macros for moduli stacks/spaces
%
\def\QCohstack{\mathcal{QC}\!\mathit{oh}}
\def\Cohstack{\mathcal{C}\!\mathit{oh}}
\def\Spacesstack{\mathcal{S}\!\mathit{paces}}
\def\Quotfunctor{\mathrm{Quot}}
\def\Hilbfunctor{\mathrm{Hilb}}
\def\Curvesstack{\mathcal{C}\!\mathit{urves}}
\def\Polarizedstack{\mathcal{P}\!\mathit{olarized}}
\def\Complexesstack{\mathcal{C}\!\mathit{omplexes}}
% \Pic is the operator that assigns to X its picard group, usage \Pic(X)
% \Picardstack_{X/B} denotes the Picard stack of X over B
% \Picardfunctor_{X/B} denotes the Picard functor of X over B
\def\Pic{\mathop{\mathrm{Pic}}\nolimits}
\def\Picardstack{\mathcal{P}\!\mathit{ic}}
\def\Picardfunctor{\mathrm{Pic}}
\def\Deformationcategory{\mathcal{D}\!\mathit{ef}}
\newenvironment{reference}{}{}
\newenvironment{slogan}{}{}
\newenvironment{history}{}{}
\usepackage{cleveref}
\crefname{originaldefinition}{original definition}{original definitions}
\Crefname{originaldefinition}{Original definition}{Original definitions}
\makeatletter
\newlabel{algebra-definition-geometrically-reduced}{{030S}{source}{}{sourcepointer-1}{}}
\newlabel{algebra-definition-geometrically-reduced@cref}{{[originaldefinition][0][]030S}{[1][0][]source}{}{}{}}
\makeatother
\crefname{originallemma}{original lemma}{original lemmas}
\Crefname{originallemma}{Original lemma}{Original lemmas}
\makeatletter
\newlabel{algebra-lemma-p-ring-map}{{0BRA}{source}{}{sourcepointer-2}{}}
\newlabel{algebra-lemma-p-ring-map@cref}{{[originallemma][0][]0BRA}{[1][0][]source}{}{}{}}
\makeatother
\crefname{originallemma}{original lemma}{original lemmas}
\Crefname{originallemma}{Original lemma}{Original lemmas}
\makeatletter
\newlabel{algebra-lemma-characterize-local-ring}{{00E9}{source}{}{sourcepointer-3}{}}
\newlabel{algebra-lemma-characterize-local-ring@cref}{{[originallemma][0][]00E9}{[1][0][]source}{}{}{}}
\makeatother
\crefname{originallemma}{original lemma}{original lemmas}
\Crefname{originallemma}{Original lemma}{Original lemmas}
\makeatletter
\newlabel{lemma-change-fields-flat}{{0C4Y}{source}{}{sourcepointer-4}{}}
\newlabel{lemma-change-fields-flat@cref}{{[originallemma][0][]0C4Y}{[1][0][]source}{}{}{}}
\makeatother
\crefname{originallemma}{original lemma}{original lemmas}
\Crefname{originallemma}{Original lemma}{Original lemmas}
\makeatletter
\newlabel{algebra-lemma-geometrically-reduced-finite-purely-inseparable-extension}{{030V}{source}{}{sourcepointer-5}{}}
\newlabel{algebra-lemma-geometrically-reduced-finite-purely-inseparable-extension@cref}{{[originallemma][0][]030V}{[1][0][]source}{}{}{}}
\makeatother
\crefname{originallemma}{original lemma}{original lemmas}
\Crefname{originallemma}{Original lemma}{Original lemmas}
\makeatletter
\newlabel{morphisms-lemma-base-change-surjective}{{01S1}{source}{}{sourcepointer-6}{}}
\newlabel{morphisms-lemma-base-change-surjective@cref}{{[originallemma][0][]01S1}{[1][0][]source}{}{}{}}
\makeatother
\title{Geometric reducedness at a point may be checked on the local ring}
\author{The Stacks Project Authors}
\date{stacks-project:varieties.tex:local-snapshot-20261011; GNU FDL 1.2 or later}
\begin{document}
\maketitle
\noindent\textit{Editorial note.} All delivered mathematical text is the author's own, copied from the pinned Stacks Project chapter file varieties.tex; no statement or proof step was rewritten, reordered or removed. The counted statement carries the author's own slogan 'Geometric reducedness can be checked on local rings', and the varieties chapter preamble already provides the project's no-op definition of the slogan environment.\par\medskip
\noindent\textit{Editorial note.} The preamble is the chapter's own retained minimal preamble (sources/stacks-varieties/retained-preamble.tex, lines 1-77), a verbatim excerpt of the native Xy-pic setup, of the native theorem declarations with mathematical macros, and of the project's own no-op definitions of the reference, slogan and history display environments; no macro was invented or redefined.\par\medskip
\noindent\textit{Editorial note.} The same-chapter reference is located at its \textbackslash{}label line inside the pinned chapter file; the five cross-chapter pointers are located in the project's own tag list staged in the snapshot (sources/stacks-chow/tags.txt), at the verbatim lines carrying their canonical tags.\par\medskip
\noindent\textit{Editorial note.} Overlap pre-check: this label occurs in no other staged chapter (run with agent-authoring/stk/card\_plan.py varieties before opening).\par\medskip
\noindent\textit{Allowed standard prerequisites (editorial).}\begin{itemize}
\item Geometrically reduced ring over a field, and its point-wise form. Algebra, Definition labelled algebra-definition-geometrically-reduced (canonical Stacks tag 030S); named in the counted statement at line 366. Supplies the algebraic notion on the right-hand side of the equivalence.
\item Local rings of a base change of schemes over a field. Same chapter, Lemma labelled lemma-change-fields-flat (source line 238, Stacks tag 0C4Y); invoked at proof lines 378-380 and again at 391-393. Supplies the description of the local rings of X\_\{k'\} as localisations of O\_\{X,x\} tensor\_k k', which is what connects the geometric condition to the ring condition in both directions.
\item The map on spectra is unchanged by a purely inseparable extension, and local rings are characterised by their spectrum. Algebra, Lemma algebra-lemma-p-ring-map (0BRA, line 374) and Lemma algebra-lemma-characterize-local-ring (00E9, line 377). Supply that Spec(O\_\{X,x\} tensor\_k k') = Spec(O\_\{X,x\}) for finite purely inseparable k'/k, so that the base change stays local with a unique point over x.
\item Surjectivity of a base change and the purely inseparable criterion. Morphisms, Lemma morphisms-lemma-base-change-surjective (01S1, line 389) and Algebra, Lemma algebra-lemma-geometrically-reduced-finite-purely-inseparable-extension (030V, line 383). Supply surjectivity of X\_\{k'\} -> X and the passage from the finite purely inseparable cases back to geometric reducedness of the ring.
\end{itemize}\medskip
\ifdefined\hypertarget\hypertarget{sourcepointer-1}{}\fi
\noindent\textit{Source reference.} Original definition 030S: sources/stacks-chow/tags.txt, source lines 3959--3959. Cross-chapter pointer named in the counted statement at line 366 for the ring-theoretic notion; the equivalence proof is retained in the counted span.\par\medskip
\ifdefined\hypertarget\hypertarget{sourcepointer-2}{}\fi
\noindent\textit{Source reference.} Original lemma 0BRA: sources/stacks-chow/tags.txt, source lines 14763--14763. Cross-chapter standard prerequisite invoked once at proof line 374 to see that a purely inseparable extension does not change the spectrum.\par\medskip
\ifdefined\hypertarget\hypertarget{sourcepointer-3}{}\fi
\noindent\textit{Source reference.} Original lemma 00E9: sources/stacks-chow/tags.txt, source lines 526--526. Cross-chapter standard prerequisite invoked once at proof lines 376-377 for localness of the tensor product ring.\par\medskip
\ifdefined\hypertarget\hypertarget{sourcepointer-4}{}\fi
\noindent\textit{Source reference.} Original lemma 0C4Y: sources/stacks-varieties/pinned-original-source.tex, source lines 238--238. Same-chapter pointer invoked twice, at proof lines 378-380 and 391-393, for the description of the local rings of a base change; it is the technical input of both directions and the surrounding argument is retained in the counted span.\par\medskip
\ifdefined\hypertarget\hypertarget{sourcepointer-5}{}\fi
\noindent\textit{Source reference.} Original lemma 030V: sources/stacks-chow/tags.txt, source lines 3962--3962. Cross-chapter standard prerequisite invoked once at proof lines 382-383 for the purely inseparable criterion used in the direction (1) => (2).\par\medskip
\ifdefined\hypertarget\hypertarget{sourcepointer-6}{}\fi
\noindent\textit{Source reference.} Original lemma 01S1: sources/stacks-chow/tags.txt, source lines 2322--2322. Cross-chapter standard prerequisite invoked once at proof lines 387-389 for surjectivity of X\_\{k'\} -> X; the localisation argument on top of it is retained in the counted span.\par\medskip
\begin{lemma}
\label{lemma-geometrically-reduced-at-point}
\begin{slogan}
Geometric reducedness can be checked on local rings.
\end{slogan}
Let $k$ be a field.
Let $X$ be a scheme over $k$.
Let $x \in X$.
The following are equivalent
\begin{enumerate}
\item $X$ is geometrically reduced at $x$, and
\item the ring $\mathcal{O}_{X, x}$ is geometrically
reduced over $k$ (see
Algebra, Definition \ref{algebra-definition-geometrically-reduced}).
\end{enumerate}
\end{lemma}

\begin{proof}
Assume (1). This in particular implies that $\mathcal{O}_{X, x}$
is reduced. Let $k'/k$ be a finite purely inseparable field
extension. Consider the ring $\mathcal{O}_{X, x} \otimes_k k'$.
By Algebra, Lemma \ref{algebra-lemma-p-ring-map}
its spectrum is the same as the spectrum of $\mathcal{O}_{X, x}$.
Hence it is a local ring also
(Algebra, Lemma \ref{algebra-lemma-characterize-local-ring}).
Therefore there is a unique point $x' \in X_{k'}$ lying over $x$
and $\mathcal{O}_{X_{k'}, x'} \cong \mathcal{O}_{X, x} \otimes_k k'$,
see Lemma \ref{lemma-change-fields-flat}.
By assumption this is a reduced ring. Hence we deduce (2) by
Algebra, Lemma
\ref{algebra-lemma-geometrically-reduced-finite-purely-inseparable-extension}.

\medskip\noindent
Assume (2). Let $k'/k$ be a field extension. Since
$\Spec(k') \to \Spec(k)$ is surjective, also
$X_{k'} \to X$ is surjective
(Morphisms, Lemma \ref{morphisms-lemma-base-change-surjective}).
Let $x' \in X_{k'}$ be any point lying over $x$.
The local ring $\mathcal{O}_{X_{k'}, x'}$
is a localization of the ring $\mathcal{O}_{X, x} \otimes_k k'$
by Lemma \ref{lemma-change-fields-flat}.
Hence it is reduced by assumption and (1) is proved.
\end{proof}

\end{document}
